\documentclass[11pt,a4paper]{article}

\usepackage{fontspec}
\setmainfont{DejaVu Serif}[
  BoldFont={DejaVu Serif Bold},
  ItalicFont={DejaVu Serif},
  ItalicFeatures={FakeSlant=0.18},
  BoldItalicFont={DejaVu Serif Bold},
  BoldItalicFeatures={FakeSlant=0.18}
]
\setsansfont{DejaVu Sans}
\usepackage{polyglossia}
\setmainlanguage{english}
\usepackage[a4paper,margin=20.5mm]{geometry}
\usepackage{microtype}
\usepackage{setspace}
\setstretch{1.035}
\usepackage{amsmath,amssymb,amsthm,mathtools,bm}
\usepackage{booktabs,longtable,array,tabularx}
\usepackage{enumitem}
\usepackage{xcolor}
\usepackage[unicode,hidelinks]{hyperref}
\usepackage[nameinlink,noabbrev]{cleveref}
\emergencystretch=3em
\allowdisplaybreaks

\hypersetup{
  pdftitle={Roller-Coaster Geometry as a Typed Observation Laboratory},
  pdfauthor={Stas, Independent Research Program},
  pdfsubject={Framed curves, specific force, rail offsets, dissipation, and observation loss},
  pdfkeywords={roller coaster, framed curve, Bishop frame, specific force, offset curve, observation geometry}
}

\definecolor{obsblue}{RGB}{37,83,138}
\definecolor{obsred}{RGB}{150,45,45}
\definecolor{lightblue}{RGB}{239,246,253}
\definecolor{lightred}{RGB}{253,242,242}

\theoremstyle{definition}
\newtheorem{definition}{Definition}[section]
\newtheorem{model}[definition]{Model}
\theoremstyle{plain}
\newtheorem{theorem}[definition]{Theorem}
\newtheorem{proposition}[definition]{Proposition}
\newtheorem{corollary}[definition]{Corollary}
\theoremstyle{remark}
\newtheorem{remark}[definition]{Remark}
\newtheorem{warning}[definition]{Boundary}

\newcommand{\R}{\mathbb{R}}
\newcommand{\dd}{\mathop{}\!d}
\newcommand{\norm}[1]{\left\lVert #1\right\rVert}
\newcommand{\SO}{\mathrm{SO}}
\newcommand{\SE}{\mathrm{SE}}
\newcommand{\Ocal}{\mathcal O}
\newcommand{\Kcal}{\mathcal K}
\newcommand{\Rcal}{\mathcal R}
\newcommand{\crossop}[1]{\left[#1\right]_\times}
\newcolumntype{L}[1]{>{\raggedright\arraybackslash}p{#1}}
\newcolumntype{Y}{>{\raggedright\arraybackslash}X}

\title{\bfseries Roller-Coaster Geometry\\
\large Framed curves, rider observations, and regular rail offsets}
\author{Stas\\
\small Independent Research Program; ORCID: 0009-0000-2294-705X}
\date{Strict GO Core reference revision v1.1 --- 28 July 2026}

\begin{document}
\maketitle

\begin{center}
\fcolorbox{obsblue}{lightblue}{%
\parbox{0.92\linewidth}{\small
\textbf{Status.} This document supersedes the bilingual working draft
v1.0. It is migrated to GO Core v0.2 and the Frames--Forces--Dissipation
Interface v0.1. It is a differential-geometric and observation
laboratory, not a track design, vehicle model, human-tolerance model,
or safety certificate.}}
\end{center}

\begin{abstract}
A roller-coaster centerline does not encode vehicle orientation,
rail-contact paths, or rider measurements. We define a track as a
unit-speed curve \(\gamma(s)\), a nonsingular adapted frame
\(R(s)\in\SO(3)\), and either a prescribed speed profile or a declared
dynamical law. With \(R^{\mathsf T}R'=\crossop{\kappa_R}\), the vehicle
angular velocity is \(\Omega^B=v\kappa_R\). The inertial acceleration
is
\[
a^I=v\frac{\dd v}{\dd s}T+v^2T',
\]
while an ideal vehicle accelerometer reads the body-frame specific
force
\[
f^B=R^{\mathsf T}(a^I-g^I).
\]
The dimensionless load vector
\(\bm G=f^B/g_0\) and its scalar norm \(G=\norm{\bm G}\) are separate
quantities. Body-frame jerk includes the frame-rotation correction and
is not merely \(R^{\mathsf T}\dot a\). Ideal rail offsets
\(\gamma_\pm=\gamma\pm(d/2)e_2\) are regular under the sufficient bound
\((d/2)\sup\norm{e_2'}<1\); global nonintersection additionally
requires a tubular-neighborhood hypothesis. The labeled length split
\(\Delta L=L_+-L_-\) changes sign under rail exchange, whereas
\(\lvert\Delta L\rvert\) is exchange invariant. Dissipation and moving
support work follow the common energy contract. All claims are
qualified by their transformation group and model status.
\end{abstract}

\tableofcontents

\section{Scope and protocol}

An observation protocol \(\lambda\) specifies the world frame,
adapted-frame convention, clock, arc-length chart, sensor axes,
sampling, calibration, stochastic/noiseless channel, estimator, and
uncertainty. The chain
\[
X_t\xrightarrow{\Rcal}\widetilde X_t
\xrightarrow{\Ocal_\lambda}Z_t
\xrightarrow{\Kcal_\lambda}\operatorname{Law}(Y_t\mid Z_t)
\xrightarrow{\widehat\theta_\lambda}\widehat\theta_t
\]
keeps the physical state, ideal sensor signal, recorded data, and
estimate distinct.

\begin{warning}[Two admissible model modes]
A prescribed pair \((\gamma,v)\) defines a kinematic benchmark. A
dynamical prediction requires mass, forces, constraints, drive,
dissipation, and initial data. The speed may not be simultaneously
prescribed and claimed as the output of an unspecified force law.
\end{warning}

\section{Track as an adapted framed curve}

Let
\[
\gamma:[0,L]\to\R^3,\qquad
\norm{\gamma'(s)}=1,
\]
be \(C^3\), and set \(T=\gamma'\). An adapted frame is
\[
R(s)=(e_1,e_2,e_3)\in\SO(3),\qquad e_1=T.
\]
It may be a Bishop/rotation-minimizing frame with an added roll
profile, or any other declared \(C^2\) adapted frame. Unlike the
Frenet frame, this definition remains meaningful at points where
\(\gamma''=0\), provided \(R\) is supplied smoothly.

\begin{definition}[Frame strain]
The spatial frame strain \(\kappa_R^B(s)\) is defined by
\begin{equation}\label{eq:frame-strain}
R^{\mathsf T}R'=\crossop{\kappa_R^B},
\qquad
\crossop{\kappa}u=\kappa\times u .
\end{equation}
It has dimension \(L^{-1}\).
\end{definition}

Curvature magnitude is
\[
\kappa(s)=\norm{T'(s)}.
\]
Torsion is defined through the Frenet frame only where
\(\kappa>0\) with sufficient regularity. The supplied roll/frame
profile remains the primary apparatus orientation.

\section{Time kinematics}

Let \(s=s(t)\) with \(v=\dot s>0\). Then
\begin{align}
\dot\gamma&=vT,\\
a^I=\ddot\gamma
&=\dot v\,T+v^2T'
=v\frac{\dd v}{\dd s}T+v^2T'.
\label{eq:acceleration}
\end{align}
Since
\[
R^{\mathsf T}\dot R
=vR^{\mathsf T}R'
=\crossop{v\kappa_R^B},
\]
the vehicle-frame angular velocity is
\begin{equation}\label{eq:angular-velocity}
\Omega^B=v\kappa_R^B.
\end{equation}

Let \(g^I\) be gravitational acceleration and let \(g_0>0\) be a
declared reference acceleration. The ideal accelerometer specific
force is
\begin{equation}\label{eq:specific-force}
f^B=R^{\mathsf T}(a^I-g^I).
\end{equation}

\begin{definition}[Vector and scalar load]
\begin{equation}\label{eq:load-definitions}
\bm G(s)=\frac{f^B(s)}{g_0},
\qquad
G(s)=\norm{\bm G(s)}.
\end{equation}
\end{definition}

\(\bm G\) is a dimensionless body-frame vector; \(G\) is a
dimensionless scalar. They cannot share one quantity identifier.

\begin{definition}[Body-frame jerk signal]
The derivative of the measured specific-force components is
\begin{equation}\label{eq:body-jerk}
j_f^B
=\frac{\dd f^B}{\dd t}
=R^{\mathsf T}(\dot a^I-\dot g^I)
-\Omega^B\times f^B.
\end{equation}
\end{definition}

The last term is required because the sensor axes rotate. The inertial
jerk \(\dot a^I\), the derivative of specific force, and the derivative
of body components are different observables.

\section{Rail offsets and their domain}

Let \(d>0\) be the ideal rail gauge and let \(e_2(s)\) be the labeled
lateral frame direction. Define
\begin{equation}\label{eq:offsets}
\gamma_\pm(s)=\gamma(s)\pm\frac d2e_2(s).
\end{equation}

\begin{proposition}[Sufficient local regularity bound]
If
\begin{equation}\label{eq:offset-bound}
\frac d2\sup_{s\in[0,L]}\norm{e_2'(s)}<1,
\end{equation}
then both offset curves in \cref{eq:offsets} are regular.
\end{proposition}

\begin{proof}
\[
\norm{\gamma_\pm'}
=\norm{T\pm(d/2)e_2'}
\ge1-(d/2)\norm{e_2'}>0.
\]
\end{proof}

The bound is sufficient, not necessary. It prevents a local offset
cusp but not a global self-intersection. To interpret the offsets as
two distinct rails, the strip map
\[
\Psi(s,u)=\gamma(s)+u e_2(s),
\qquad |u|\le d/2,
\]
must additionally be an embedding (or be restricted to a declared
tubular neighborhood).

The rail lengths are
\[
L_\pm=\int_0^L\norm{\gamma_\pm'(s)}\,\dd s.
\]
For a labeled frame define
\[
\Delta L_{\rm or}=L_+-L_-,
\qquad
\Delta L_{\rm abs}=\lvert L_+-L_-\rvert.
\]

\begin{proposition}[Transformation law of the split]
Under a static Euclidean isometry applied to the entire framed track,
\(L_\pm\), \(\Delta L_{\rm or}\), and
\(\Delta L_{\rm abs}\) are unchanged. Under the label exchange
\(+\leftrightarrow-\),
\[
\Delta L_{\rm or}\mapsto-\Delta L_{\rm or},
\qquad
\Delta L_{\rm abs}\mapsto\Delta L_{\rm abs}.
\]
\end{proposition}

Thus the signed split is a covariant of an oriented/labeled rail
system, not an unqualified invariant.

\section{Rolling count and contact qualification}

For constant effective wheel radius \(r_w>0\), pure rolling without
slip along rail \(\gamma_\pm\) gives the ideal revolution count
\[
N_\pm^{\rm ideal}=\frac{L_\pm}{2\pi r_w}.
\]
This formula is dimensionless and parameterization independent. It is
not an actual wheel count if slip, flange contact, compliance, wear,
variable rolling radius, or loss of contact occurs. Those effects
require a contact model and a wheel-angle observation channel.

\section{Force and dissipation layer}

Let \(F_d^I\) be the total dissipative spatial force on a car at a
declared application point and \(v_{\rm rel}^I\) its velocity relative
to the corresponding support or medium. The common sign convention is
\[
P_{\rm diss}=-F_d^I\cdot v_{\rm rel}^I\ge0.
\]
For a stationary ideal rail constraint, the normal reaction performs
no work along an admissible motion. With generalized Lagrangian energy
\(E_L\),
\[
\dot E_L
=P_{\rm drive}-P_{\rm diss}+P_{\rm con}-\partial_tL.
\]
For a static autonomous track with ideal normal constraints,
\(P_{\rm con}=0\). A launch system, moving track element, or active
brake is not covered by that simplification.

\begin{warning}[No force from geometry alone]
Curvature determines the kinematic acceleration required by a chosen
speed. It does not by itself determine wheel loads, structural stress,
or loss of contact. Those require a mass/inertia model and solved
constraint reactions.
\end{warning}

\section{Observation loss}

Consider the scalar load observation
\[
\Ocal_G:(\gamma,R,v)\longmapsto
G(s)=\frac{\norm{a^I(s)-g^I}}{g_0}.
\]

\begin{proposition}[Scalar load is not a complete framed-track
observation]
Fix a centerline \(\gamma\) and speed \(v\). Any two smooth adapted
frames \(R_1,R_2\) with the same tangent but distinct roll profiles
produce the same scalar load \(G\). Hence \(\Ocal_G\) is not injective
on framed tracks.
\end{proposition}

\begin{proof}
Equation \cref{eq:acceleration} depends on \(\gamma,v\), not on the
choice of lateral/up axes. Orthogonality gives
\[
\norm{R_i^{\mathsf T}(a^I-g^I)}
=\norm{a^I-g^I}
\]
for both frames, although their body-component vectors generally
differ.
\end{proof}

The full vector \(\bm G\), angular velocity, visual flow, contact
signals, and known speed add information but are still protocol
outputs, not automatically a complete reconstruction of the track.

\section{Transformation audit}

\begin{longtable}{L{34mm}L{48mm}L{68mm}}
\toprule
Object & Transformation group & Status \\
\midrule
\endhead
Arc length and curvature magnitude & static Euclidean isometries &
invariant \\
Frame components \(\bm G,\Omega^B\) & declared apparatus-axis change &
covariant vectors \\
Scalar load \(G\) & apparatus-axis orthogonal change & invariant \\
Signed rail split \(\Delta L_{\rm or}\) & rail-label exchange &
sign-covariant \\
Absolute split \(\Delta L_{\rm abs}\) & rail-label exchange &
invariant \\
Ideal wheel count & curve reparameterization & invariant under
orientation-preserving regular reparameterization \\
Body jerk components & time-dependent apparatus frame & require the
correction in \cref{eq:body-jerk} \\
\bottomrule
\end{longtable}

\section{Reproducibility and claim boundary}

A numerical benchmark must report:
\begin{enumerate}[label=(\roman*)]
  \item \(\gamma,R,v\) and their differentiability;
  \item sampling and the derivative/smoothing method;
  \item the bound \cref{eq:offset-bound} and global strip validation;
  \item gravity, sensor axes, and force/contact status;
  \item numerical convergence.
\end{enumerate}
The v1.0 synthetic numbers are not retained as engineering results:
the layout supplied no clearance, structural, contact, or convergence
certificate.

This document does not claim:
\begin{itemize}
  \item a realizable or safe track;
  \item a multibody wheel--rail solution;
  \item structural, fatigue, restraint, or human-tolerance validation;
  \item recovery of a track from one rider signal;
  \item novelty of the standard curve and accelerometer formulas.
\end{itemize}

\begin{thebibliography}{9}

\bibitem{bishop}
R.~L. Bishop, ``There is more than one way to frame a curve,''
\emph{American Mathematical Monthly} 82 (1975), 246--251.
\href{https://doi.org/10.2307/2319846}{doi:10.2307/2319846}.

\bibitem{docarmo}
M.~P. do Carmo, \emph{Differential Geometry of Curves and Surfaces},
Prentice-Hall, 1976.

\bibitem{brogliato}
B. Brogliato, \emph{Nonsmooth Mechanics: Models, Dynamics and
Control}, 3rd ed., Springer, 2016.
\href{https://doi.org/10.1007/978-3-319-28664-8}
{doi:10.1007/978-3-319-28664-8}.

\end{thebibliography}

\end{document}
